\section{Auxiliary results}
\label{app:auxiliary}

This appendix collects the classical finite facts invoked by the construction
and the upper-bound proof, followed by the supplementary row-selector
invariance result (\cref{app:prop:row-selector-invariance}).
The elementary finite facts are deterministic algebraic or graph
statements and introduce no additional model assumptions.

\subsection{Finite Hellinger-affinity calculus}
\label{app:auxiliary-affinity}

\textcolor{red}{The textbook cited below uses $H^2(p,q)=2\sqHellinger(p,q)$
(\href{https://people.lids.mit.edu/yp/homepage/data/itbook-export.pdf}{author-hosted PDF}).
We include finite proofs for completeness.}

\begin{lemma}[Product affinity and TV comparison; \textcolor{red}{{\citealp[Eqs.~(7.22) and (7.26)]{polyanskiy2025information}}}]
\label{app:lem:finite-hellinger-calculus}
Let $J$ be finite.  For every $j\in J$, let $p_j,q_j$ be probability masses
on a finite set $\mathcal X_j$.  Then
\begin{equation}
 \HellingerAffinity\!\left(\bigotimes_{j\in J}p_j,
                           \bigotimes_{j\in J}q_j\right)
 =\prod_{j\in J}\HellingerAffinity(p_j,q_j),
 \label{app:eq:aux-product-affinity}
\end{equation}
and consequently
\begin{equation}
 \sqHellinger\!\left(\bigotimes_{j\in J}p_j,
                      \bigotimes_{j\in J}q_j\right)
 \le\sum_{j\in J}\sqHellinger(p_j,q_j).
 \label{app:eq:aux-product-hellinger}
\end{equation}
For any two probability masses $p,q$ on the same finite set,
\begin{equation}
 \sqHellinger(p,q)
 \le\dTV(p,q)
 \le\sqrt{2\sqHellinger(p,q)}.
 \label{app:eq:aux-tv-hellinger}
\end{equation}
\end{lemma}

\begin{proof}\color{red}
\emph{1. Products.}
For $f_j(x):=\sqrt{p_j(x)q_j(x)}$, distributivity gives
$\HellingerAffinity(\bigotimes_jp_j,\bigotimes_jq_j)
=\sum_{(x_j)_j\in\prod_j\mathcal X_j}\prod_j f_j(x_j)
=\prod_j(\sum_{x\in\mathcal X_j}f_j(x))
=\prod_j\HellingerAffinity(p_j,q_j)$.
Put $u_j:=\sqHellinger(p_j,q_j)\in[0,1]$, enumerate $J$, and telescope:
$1-\prod_j(1-u_j)
=\sum_j[\prod_{k<j}(1-u_k)-\prod_{k\le j}(1-u_k)]
=\sum_j u_j\prod_{k<j}(1-u_k)\le\sum_j u_j$.
Since $\sqHellinger=1-\HellingerAffinity$, this and
\eqref{app:eq:aux-product-affinity} give \eqref{app:eq:aux-product-hellinger}.

\emph{2. TV comparison.}
The identity $\min\{a,b\}=(a+b-|a-b|)/2$ gives
$\sum_x\min\{p(x),q(x)\}
=(2-\sum_x|p(x)-q(x)|)/2=1-\dTV(p,q)$.
Since \(\min\{p(x),q(x)\}\le\sqrt{p(x)q(x)}\), subtracting the summed
inequality from one gives
\(\sqHellinger(p,q)\le\dTV(p,q)\).  For the other direction,
Cauchy--Schwarz yields
\begin{align*}
 2\dTV(p,q)
 &=\sum_x|\sqrt{p(x)}-\sqrt{q(x)}|
          (\sqrt{p(x)}+\sqrt{q(x)})\\
 &\le
 \left(\sum_x(\sqrt{p(x)}-\sqrt{q(x)})^2\right)^{1/2}
 \left(\sum_x(\sqrt{p(x)}+\sqrt{q(x)})^2\right)^{1/2}\\
 &\le 2\sqrt{2\sqHellinger(p,q)}.
\end{align*}
Here the first squared sum is
$\sum_x(\sqrt{p(x)}-\sqrt{q(x)})^2
=2-2\HellingerAffinity(p,q)=2\sqHellinger(p,q)$, and the second is
$\sum_x(\sqrt{p(x)}+\sqrt{q(x)})^2
=\sum_xp(x)+\sum_xq(x)+2\sum_x\sqrt{p(x)q(x)}
=2+2\HellingerAffinity(p,q)\le4$.
The last inequality uses Cauchy--Schwarz:
$\HellingerAffinity(p,q)\le\sqrt{\sum_xp(x)\sum_xq(x)}=1$.
Dividing by two proves \eqref{app:eq:aux-tv-hellinger}.
\end{proof}

\begin{lemma}[KL-to-Hellinger comparison; \textcolor{red}{{\citealp[Eq.~(7.33)]{polyanskiy2025information}}}]
\label{app:lem:kl-hellinger}
For strictly positive probability masses $p,q$ on a finite set,
\begin{equation}
 \sqHellinger(p,q)
 \le 1-\exp\{-\KLDivergence{p}{q}/2\}
 \le\frac12\KLDivergence{p}{q}.
 \label{app:eq:aux-kl-hellinger}
\end{equation}
The same inequalities hold with the KL arguments reversed.
\end{lemma}

\begin{proof}\color{red}
This is the order-$1/2$ R\'enyi/KL comparison in our Hellinger
normalization (see also \citealp{vanerven2014renyi}).
Jensen's inequality for the convex exponential gives
$\HellingerAffinity(p,q)
=\mathbb E_{X\sim p}\exp\{-\tfrac12\log\frac{p(X)}{q(X)}\}
\ge\exp\{-\KLDivergence{p}{q}/2\}$.
Subtract from one and use $1-e^{-x}\le x$ for $x\ge0$.
The reversed statement follows by exchanging $p,q$ and using the symmetry
of Hellinger distance.
\end{proof}

\begin{remark}[No upper KL bound from Hellinger alone]
\label{app:rem:no-hellinger-kl-converse}
Strict positivity without a uniform lower bound on the masses does not
give a converse upper bound on KL.  For $0<u<1/4$, consider
\[
 p_u=(1-u,u),\qquad
 q_u=(1-u e^{-1/u^2},u e^{-1/u^2}).
\]
Both laws are strictly positive, and
\[
 \sqHellinger(p_u,q_u)\le\dTV(p_u,q_u)
 =u(1-e^{-1/u^2})\le u\longrightarrow0.
\]
Nevertheless,
\[
 \KLDivergence{p_u}{q_u}
 =\frac1u+(1-u)\log\frac{1-u}{1-u e^{-1/u^2}}
 \ge\frac1u-u\longrightarrow\infty,
\]
For the last inequality, $1-u e^{-1/u^2}\le1$ implies
$\log((1-u)/(1-u e^{-1/u^2}))\ge\log(1-u)$, while
\[
 -\log(1-u)=\int_0^u\frac{dt}{1-t}
 \le\frac{u}{1-u}
 \quad\Longrightarrow\quad (1-u)\log(1-u)\ge-u.
\]
Swapping the two laws proves the corresponding failure for reverse KL.
\end{remark}

\subsection{Elementary forest facts}
\label{app:auxiliary-forest}

\begin{lemma}[Forest degree sum and tree centroid]
\label{app:lem:forest-basics}
Let $H=(W,E_H)$ be a finite nonempty forest with $c(H)$ connected
components.  Then
\begin{equation}
 |E_H|=|W|-c(H),
 \qquad
 \sum_{v\in W}d_H(v)=2|E_H|<2|W|.
 \label{app:eq:aux-forest-degree-sum}
\end{equation}
Moreover, every finite nonempty tree $T$ has a vertex $v$ such that every
component of $T\setminus\{v\}$ has at most $|V(T)|/2$ vertices.
\end{lemma}

\begin{proof}
Each tree component with $n_\ell$ vertices has $n_\ell-1$ edges.  Summing
over the components gives the first identity; the degree-sum identity then
gives the remaining claims in \eqref{app:eq:aux-forest-degree-sum}.

For the centroid claim, put $n:=|V(T)|$ and choose $v$ minimizing the largest component size
after deleting $v$.  If a component $S$ of $T\setminus\{v\}$ had more than
$|V(T)|/2$ vertices, let $w$ be the neighbor of $v$ in $S$.  After deleting
$w$, the component containing $v$ has $|V(T)|-|S|<|V(T)|/2$ vertices, while
every other component is a proper subset of $S$ and hence has at most
$|S|-1$ vertices.  Thus the largest component after deleting $w$ has size
at most
\[
 \max\{n-|S|,\ |S|-1\}<|S|,
\]
whereas the largest component after deleting $v$ has size $|S|$
(all other components together have $n-1-|S|<|S|$ vertices).
This contradicts the choice of $v$, proving the centroid claim.
\end{proof}

\subsection{Public row-selector invariance}

\paragraph{Intuition.}
The structural proof uses two selector properties: exact candidate rows
of size at most $\DegreeCutoff$ are preserved, and every returned row
has size at most $\DegreeCutoff$.  Consequently the same successful-path and all-path arguments apply
to every public selector with this contract.  The statistical comparisons
are then constructed for that selector's own decision rule.

\begin{proposition}[Invariance under public row selection]
\label{app:prop:row-selector-invariance}
Replace the standard top-vote selector by any rule in
Definition~\ref{app:def:row-selector}, keeping all other screen parameters,
fresh-color rules, and decision-rule caps unchanged.
Under the finite hypotheses of Theorem~\ref{app:thm:finite-upper},
$\SelectedPackedSampler$ has the same low-degree screen contract,
successful-path contraction and stopping guarantees, all-path resource
bounds, and expected-TV bound $\OracleRadius+\ScreenFailureBudget$.
Under the additional oracle hypotheses of
Theorem~\ref{app:thm:finite-kl-upper}, its expected forward KL is at most
$N\KLRowBudget+N\log\VocabSize\,\ScreenFailureBudget$.
Consequently both main-text benchmark guarantees retain their original
calibrations and rates.  This statement compares guarantees, not the
realized paths or errors of two selectors.
\end{proposition}

\begin{proof}

Recall the selector contract of \cref{app:def:row-selector}:
\[
 \begin{gathered}
 \RowSelector(j,C,\DegreeCutoff;\ScreenTranscript)\subseteq C,
 \qquad
 |\RowSelector(j,C,\DegreeCutoff;\ScreenTranscript)|\le\DegreeCutoff,\\
 |C|\le\DegreeCutoff
 \quad\Longrightarrow\quad
 \RowSelector(j,C,\DegreeCutoff;\ScreenTranscript)=C.
 \end{gathered}
\]
The selector is a deterministic function of the available public transcript,
makes no submission or commit, and does not read future randomness.
The dependencies of the proof are as follows.

\begin{center}
\small
\begin{tabularx}{\linewidth}{@{}>{\raggedright\arraybackslash}p{.28\linewidth}
 >{\raggedright\arraybackslash}p{.30\linewidth}
 >{\raggedright\arraybackslash}X@{}}
\toprule
Property retained & Argument using it & Guarantee retained\\
\midrule
Identity on nonoverflowing rows
 & \AppLowDegreeScreenRef & Low-degree exactness\\
Row-size cap and low-degree exactness
 & \Cref{app:lem:successful-path-structure,app:lem:upper-resources}
 & Contraction, stopping, resource bounds\\
Fresh colors conditional on the past
 & \Cref{app:lem:random-screen-success} & Failure budget\\
This selector's own safe rule
 & \Cref{app:lem:exact-hybrid,app:lem:adaptive-hellinger-upper,app:lem:safe-controller-comparison}
 & Expected TV bound\\
This selector's own joint reference
 & \Cref{app:lem:adaptive-joint-kl,app:thm:finite-kl-upper}
 & Expected forward-KL bound\\
\bottomrule
\end{tabularx}
\end{center}

\emph{1. Preserve the low-degree screen contract.}
On a successful screen, the proof of
\AppLowDegreeScreenRef{} identifies the candidate row with
the true neighborhood before applying the selector in \AlgoBlockRef{aggregate}.
Its size is at most
$\DegreeCutoff$, so the nonoverflow identity preserves it.  Every other
returned row still has size at most $\DegreeCutoff$; no claim that it
contains only true neighbors is needed.

\emph{2. Transfer the structural and resource bounds.}
These two properties are exactly those used in
\eqref{app:eq:survivor-high-neighbors} and
\eqref{app:eq:total-claim-bound}.  Thus the contraction, terminal exactness,
and successful-path round count above apply to the decision rule with
the chosen selector.
In particular, neither cycle repair nor the round-cap fallback activates
on a successful path.
The selector adds no submission or commit, and the unchanged caps bound
resources on every path, including failed screens.


\emph{3. Retain the conditional screen-failure budget.}
Conditioning on this decision rule's realized past, including previous
selector outputs, still leaves the next colors drawn in \AlgoBlockRef{draft}
independent and uniform.  Lemma~\ref{app:lem:random-screen-success} therefore gives the
same failure budget without a union bound over selectors or histories.

\emph{4. Use the chosen selector's own safe decision rule for TV.}
For TV, construct the safe decision rule of
Lemma~\ref{app:lem:safe-controller-comparison} using this same selector.
The common-path coupling compares $\SelectedPackedSampler$ with its own
safe decision rule using the same selector; it does not compare different
selectors.  Structural safety and the generic adaptive Hellinger bound
give $\OracleRadius$, and a failed screen contributes at most
$\ScreenFailureBudget$.


\emph{5. Use the chosen selector's own joint-reference process for KL.}
For KL, instead construct the joint-reference process of
Appendix~\ref{app:kl-joint-reference} for the decision rule with the chosen
selector.  Write $\ScreenFailureEvent$ for the event that a reached
screen fails in this selector's joint-reference process.
On its reference paths, the costs in
\eqref{app:eq:adaptive-joint-kl} satisfy
\begin{align*}
 \sum_t\sum_{j\in B(\mathfrak h_t)}
  \KLDivergence{\ExactRow{y(\mathfrak h_t)}{j}}
               {\OracleRow{y(\mathfrak h_t)}{j}}
 &\le\KLRowBudget\sum_t|B(\mathfrak h_t)|=N\KLRowBudget,\\
 \sum_t\mathrm{TC}(\JointBatchLaw{\mathfrak h_t})
 &\le N\log\VocabSize\,\mathbf 1_{\ScreenFailureEvent}.
\end{align*}
The second line is zero on a successful reference path by the structural
argument; on any path, the sum is at most $N\log\VocabSize$.
Conditioning before each fresh color block in this same reference process
gives \eqref{app:eq:joint-reference-screen-failure}.  Therefore
\begin{align*}
 \mathbb E_W\KLDivergence{\TargetLaw}{\SamplerOutputLaw{W}{\FrozenOracle}}
 &\le N\KLRowBudget+
       N\log\VocabSize\,\mathbb E_W\JointReferenceLaw{W}(\ScreenFailureEvent)\\
 &\le N\KLRowBudget+N\log\VocabSize\,\ScreenFailureBudget.
\end{align*}
All output and reference laws here belong to the chosen selector.
No comparison with the standard top-vote selector, or conversion from the
TV coupling, is used.
\end{proof}

\begin{proposalcontext}{Interpreting selector invariance}
The resource guarantees do not bound a selector's local arithmetic or
memory costs.
\end{proposalcontext}

\ifincludeexponents
\input{src/sections/proofs/lower_oracle_floor}
\fi
